% Part: methods
% Chapter: proofs
% Section: example-1

\documentclass[../../../include/open-logic-section]{subfiles}

\begin{document}

\olfileid{mth}{prf}{ex1}

\olsection{يوه بېلګه}

زموږ لومړۍ بېلګه د سټونو د اتحاد او تقاطع
په اړه لاندې ساده خبره ده۔ دا د تعريفونو
پرانيستل، د عطف او عمومي دعوو ثابتول،
او د حالتونو له مخې ثبوت روښانوي۔

\begin{prop}
د هر ډول سټونو $A$، $B$، او $C$ دپاره،
$A \cup (B \cap C) = (A \cup B)
\cap (A \cup C)$۔
\end{prop}

راځئ ثابته يې کړو!{}

\begin{proof}
غواړو وښيو چې د هر ډول سټونو $A$، $B$، او
$C$ دپاره $A \cup (B \cap
C) = (A \cup B) \cap (A \cup C)$۔
\begin{quote}
لومړے د بيان د «$=$» تعريف پرانيزو۔
په ياد ولرئ چې د سټونو يو شان کېدل
ثابتول دا معنا لري چې دواړه يو شان
!!{element}s لري۔ يعنې د $A \cup
(B \cap C)$ ټول !!{element}s
د $(A \cup B) \cap (A \cup C)$
!!{element}s هم دي، او برعکس۔
«برعکس» معنا دا ده چې د $(A
\cup B) \cap (A \cup C)$ هر
!!{element} بايد د $A \cup (B \cap
C)$ !!a{element} هم وي۔
نو د تعريف له پرانيستلو ښکاري
چې يو عطف ثابتول پکار دي۔
دا وليکو:
\end{quote}
د تعريف له مخې،
$A \cup (B \cap C) = (A \cup B) \cap (A \cup C)$
هغه وخت او يوازې هغه وخت وي چې د
$A \cup (B \cap C)$ هر !!{element}
د $(A
\cup B) \cap (A \cup C)$
!!a{element} هم وي، او د
$(A \cup B) \cap (A
\cup C)$ هر !!{element} د
$A \cup (B \cap C)$
!!a{element} وي۔
\begin{quote}
څرنګه چې دا عطف دے، د دواړو غړو
ثبوت جلا پکار دے۔ له لومړي څخه
پيلوو: ثابتوو چې د
$A \cup (B \cap C)$ هر
!!{element} د $(A
\cup B) \cap (A \cup C)$
!!a{element} هم دے۔

دا عمومي دعوه ده؛ نو د
$A \cup (B \cap C)$ يو هر ډول
!!{element} نيسو او ښيو چې
د $(A \cup B) \cap (A \cup C)$
!!a{element} هم دے۔ د دغه
هر ډول !!{element} دپاره يو
متغير، لکه~$z$، ټاکو۔ ثبوت
مو داسې دوام کوي:
\end{quote}
لومړے ثابتوو چې د
$A \cup (B \cap C)$ هر
!!{element} د
$(A \cup B) \cap (A \cup C)$
!!a{element} هم دے۔ $z \in A \cup (B
\cap C)$ دې وي۔ بايد وښيو چې
$z \in (A \cup B) \cap (A \cup C)$۔
\begin{quote}
اوس د $\cup$ او~$\cap$ تعريفونه
پرانيزو۔ د بېلګې په توګه، د
$\cup$ تعريف دا دے:
$A \cup B = \Setabs{z}{z \in A
  \text{ or } z \in B}$۔ کله چې
تعريف په «$A \cup (B
\cap C)$» کاروو، د تعريف د
«$B$» رول اوس «$B \cap C$»
لوبوي؛ نو
$A \cup (B \cap C) = \Setabs{z}{z \in A \text{ or }
  z \in B \cap C}$۔
ځکه زموږ د $z \in A \cup (B \cap C)$
فرض دا معنا لري:
$z \in \Setabs{z}{z \in A \text{ or } z \in B \cap
  C}$۔ او $z \in \Setabs{z}{\dots z\dots}$
هغه وخت او يوازې هغه وخت وي چې
\dots $z$ \dots؛ يعنې دلته
$z \in A$ يا $z \in B \cap C$۔
\end{quote}
د $\cup$ له تعريفه،
$z \in A$ يا $z \in B \cap C$۔
\begin{quote}
دا فصل دے؛ نو د حالتونو له مخې
ثبوت ګټور دے۔ دواړه حالتونه اخلو
او په هر يوه کښې غوښتل شوې پايله،
يعنې «$z \in (A \cup B) \cap (A \cup
C)$»، ښيو۔
\end{quote}
حالت ۱: فرض کړئ $z \in A$۔
\begin{quote}
له خپلو فرضونو څخه نور ډېر څه
په لاس کښې نه لرو۔ نو هغه څه
وګورو چې په پايله کښې ترې کار
اخيستلے شو۔ غواړو وښيو چې
$z \in (A \cup B) \cap (A \cup C)$۔
د $\cap$ له تعريفه، د
$z \in (A \cup B) \cap (A \cup C)$
د ښودلو دپاره بايد وښيو چې
دا هم په $(A \cup B)$ او هم
په $(A \cup C)$ کښې دے۔
خو $z
\in A \cup B$ هغه وخت او يوازې
هغه وخت وي چې $z \in A$ يا
$z \in B$ وي، او د حالت~۱
له فرضه $z \in A$ راسره شته۔
په همداسې استدلال، $C$ د $B$
پر ځاے راوړلو سره، $z \in A \cup C$
راځي۔ دا استدلال مو په معکوس
لوري وکړ؛ اوس يې د ثبوت د
اړين لوري په ترتيب ليکو۔
\end{quote}
څرنګه چې $z \in A$، نو
$z \in A$ يا $z \in B$؛
له دې امله د~$\cup$ له تعريفه
$z \in A \cup B$۔ همدارنګه
$z \in A \cup C$۔ دا د
$z \in (A \cup B) \cap (A \cup C)$
معنا لري، د~$\cap$ له تعريفه۔
\begin{quote}
دا د حالتونو له مخې د ثبوت
لومړے حالت بشپړوي۔ اوس په
دويم حالت کښې، چې
$z \in B \cap C$ دے، هم
غوښتل شوې پايله اخلو۔
\end{quote}
حالت ۲: فرض کړئ $z \in B \cap C$۔
\begin{quote}
بيا د دوو سټونو له تقاطع
سره کار لرو۔ د~$\cap$
تعريف کاروو:
\end{quote}
څرنګه چې $z \in B \cap C$،
$z$ بايد هم د $B$ او هم د $C$
!!a{element} وي، د~$\cap$
له تعريفه۔
\begin{quote}
بيا خپلې پايلې ته ګورو۔
بايد وښيو چې $z$ هم
په $(A \cup B)$ او هم
په $(A \cup C)$ کښې
دے۔ حل بيا سمدستي راځي۔
\end{quote}
څرنګه چې $z \in B$،
$z \in (A \cup B)$۔
څرنګه چې $z \in C$،
$z \in (A
\cup C)$ هم۔ نو
$z \in (A \cup B) \cap (A \cup C)$۔
\begin{quote}
دلته مو د $\cup$ او $\cap$
تعريفونه بيا وکارول؛ خو ځکه
چې مخکې مو ياد کړي او ښودلې
مو ده چې که $z$ د دوو سټونو په يوه
کښې غړيتوب د هغو په اتحاد
کښې غړيتوب ورکوي، اړتيا نشته
چې دا ګامونه بيا په هماغه
تفصيل ووايو۔

د حالتونو له مخې د ثبوت
دويم حالت مو هم بشپړ کړ؛
اوس لومړۍ پايله ويلے شو۔
\end{quote}
نو که $z \in A \cup (B \cap C)$،
بيا $z \in (A \cup B) \cap (A \cup C)$۔
\begin{quote}
اوس يوازې بل لوری ښيو:
د $(A \cup B) \cap (A \cup C)$
هر !!{element} د
$A \cup (B \cap
C)$ !!a{element} هم دے۔
لکه مخکې، دا عمومي دعوه
د لومړي سټ د يوه هر ډول
!!{element} په فرضولو
او په دويم سټ کښې يې
د غړيتوب په ښودلو ثابتوو۔
راځئ دا ګام وليکو۔
\end{quote}
اوس فرض کړئ
$z \in (A \cup B) \cap (A \cup C)$۔
غواړو وښيو چې
$z \in A \cup (B \cap C)$۔
\begin{quote}
اوس زموږ فرض
$z \in (A \cup B) \cap (A
\cup C)$ دے۔ هيله ده
دا خبره به مغشوشوونکې
نه وي چې دلته هم د ثبوت
د لومړۍ برخې هماغه~$z$
کاروو۔ کله چې هغه برخه
بشپړه شوه، هلته شوي
فرضونه نور نه پلي کېږي؛
نو اوس د $z$ په اړه
نوي فرضونه کولے شو۔
که دا درته مغشوش وي،
په لاندې برخه کښې $z$
په بل متغير بدل کړئ۔

پوهېږو چې $z$ هم په
$A \cup B$ او هم په
$A \cup C$ کښې دے،
د~$\cap$ له تعريفه۔
د $\cup$ له تعريفه
دا نور هم پرانيستلے شو:
يا $z \in A$ يا
$z \in B$؛ او هم
يا $z \in A$ يا $z \in
C$۔ دا بيا د حالتونو له مخې
ثبوت ښکاري، خو «او» يې
يو څه مغشوشوي۔ ښايي فکر
وکړئ چې درې امکانونه دي:
$z$ يا په $A$، يا په
$B$، يا په $C$ کښې دے۔
خو دا تېروتنه ده۔ بايد
پام وکړو او هر فصل
جلا وڅېړو۔
\end{quote}
د $\cap$ له تعريفه،
$z \in A \cup B$ او
$z \in A \cup C$۔
د $\cup$ له تعريفه،
$z \in A$ يا $z \in B$۔
حالتونه بېلوو۔
\begin{quote}
ځکه چې اوس په لومړي فصل
تمرکز کوو، دويم فصل
(چې د $A \cup C$ له
پرانيستلو راځي) مو لا
نه دے راوړے؛ په حقيقت
کښې لا ورته اړتيا هم
نشته۔ لومړے حالت
$z \in A$ دے، او د يوه
سټ !!a{element} د هماغه
سټ او هر بل سټ د اتحاد
!!a{element} هم وي۔
نو لومړے حالت اسان دے:
\end{quote}
حالت ۱: فرض کړئ
$z \in A$۔ له دې څخه
$z \in A \cup (B \cap
C)$ راځي۔
\begin{quote}
اوس دويم حالت،
$z \in B$، اخلو۔
دلته دويم $\cup$
پرانيزو او يو بل
د حالتونو له مخې
ثبوت کوو:
\end{quote}
حالت ۲: فرض کړئ
$z \in B$۔ څرنګه چې
$z \in A \cup C$،
يا $z \in A$ يا
$z \in C$۔ حالتونه
نور هم بېلوو:

حالت ۲الف: $z \in A$۔
نو بيا هم $z \in A \cup (B \cap C)$۔
\begin{quote}
دا يو څه عجيبه و: په دې
حالت کښې مو د~$z \in B$
فرض په رښتيا نه و پکار؛
خو ستونزه نشته۔
\end{quote}
حالت ۲ب: $z \in C$۔
نو $z \in B$ او
$z \in C$، ځکه
$z \in B \cap C$؛
او له دې څخه
$z \in A \cup (B \cap C)$۔
\begin{quote}
دا د حالتونو دواړه
ثبوتونه بشپړوي؛ نو
دويمه نيمايي هم
ثابته شوه۔
\end{quote}
نو که
$z \in (A \cup B) \cap (A \cup C)$،
بيا $z \in A \cup (B \cap C)$۔
\end{proof}

\end{document}
